\chapter{Fixed-Branch Approximate Likelihoods and Same-Scalar Gradients}
\label{ch:bf-highdim-fixed-branch-same-scalar}
\label{ch:bf-highdim-hmc-research}
\label{eq:bf-hd-hmc-potential}
\label{prop:bf-hd-hmc-gradient-contract}

\section{Why Fixed-Branch Discipline Is Shared Across The Two Lanes}
\label{sec:bf-hd-fixed-branch-purpose}

Both method chapters in the restarted high-dimensional block face the same
structural problem.  The exact Bayesian target is a filtering law and normalizer
sequence, but the implemented approximation uses discrete numerical structure
that may change with parameters: cloud level, merge policy, factor branch,
coordinate domains, ranks, pivots, grids, fitting points, or map choices.  If
those structural choices move while the parameter moves, then the reported
derivative is not automatically the derivative of one declared scalar.

This chapter therefore treats fixed-branch discipline as a cross-method target
contract, not as an appendix to one particular algorithm.

\section{A Canonical Branch Record}
\label{sec:bf-hd-branch-record}

A branch record should be read as a three-ledger object,
\begin{equation}
  \mathcal H(B)=\{\mathcal H_{\rm fixed},\mathcal H_{\rm diff},\mathcal H_{\rm diag}\},
  \label{eq:bf-hd-branch-ledger}
\end{equation}
where:
\begin{itemize}[leftmargin=0.28in]
  \item \(\mathcal H_{\rm fixed}\) contains the structural choices that define the
  scalar on the saved branch,
  \item \(\mathcal H_{\rm diff}(\theta)\) contains the parameter-dependent numerical
  values recomputed on that branch,
  \item \(\mathcal H_{\rm diag}\) contains residuals, condition numbers, rank
  vectors, branch-failure markers, and other diagnostic quantities.
\end{itemize}
For the sparse-grid lane, the fixed ledger includes the cloud, one-dimensional
rule family, merge/order policy, factor family, thresholds, preprocessing
policy, and initial-condition policy.  For the retained-object lane, the fixed
ledger includes coordinate domains, basis families, fitting points, weights,
ranks, sweep order, shifts, stabilization choices, retained-object construction
policy, and the defensive mass / reference-density conventions that define the
same branch.

Two evaluations use the same branch if and only if their fixed ledgers match:
\begin{equation}
  B(\theta_a)=B(\theta_b)
  \quad\Longleftrightarrow\quad
  \mathcal H_{\rm fixed}(\theta_a)=\mathcal H_{\rm fixed}(\theta_b).
  \label{eq:bf-hd-branch-equality}
\end{equation}
The differentiable ledger is allowed to change with the parameter; otherwise the
same-scalar derivative would not be testing the actual fixed branch.

\section{Fixed-Branch Approximate Likelihoods Across The Three Operational Layers}
\label{sec:bf-hd-fixed-branch-likelihoods}

The restarted high-dimensional part now separates the operational burdens that
the old compressed block used to mix together.  The deterministic sparse-grid
lane exports a declared fixed-cloud scalar from
Chapter~\ref{ch:bf-highdim-fixed-cloud-sgqf-validation}.  The retained-object
TT/KR lane exports a declared fixed-branch scalar from
Chapter~\ref{ch:bf-highdim-squared-tt-fixed-branch-likelihoods}.  The present
chapter sits above those lane-specific constructions and states the common
same-scalar derivative logic they must both obey.

The three operational layers are:
\begin{equation}
\begin{aligned}
  \text{exact target:}&\quad \ell_T(\theta)=\sum_{t=1}^T\log Z_t(\theta),\\
  \text{declared approximate scalar:}&\quad \widehat\ell_T(\theta;B),\\
  \text{same-scalar derivative requirement:}&\quad
  \nabla_\theta \widehat\ell_T(\theta;B)
  \text{ differentiates the same declared } \widehat\ell_T.
\end{aligned}
  \label{eq:bf-hd-fixed-branch-three-layer}
\end{equation}
For the sparse-grid lane, \(B\) contains the fixed cloud, merge/order policy,
factor branch, and lane-specific veto rules.  For the retained-object lane,
\(B\) contains the fixed domains, fitting points, ranks, sweeps, shifts,
stabilization choices, and stored-field conventions that define the TT/KR scalar.
The chapter therefore treats fixed-branch discipline as a cross-method target
contract, not as an appendix to one particular algorithm.

For the actual-SV family, the ordering of obligations is now explicit.  First
one declares the target-correct scalar.  For the transformed actual-SV problem
that scalar is the cumulative likelihood built from the exact transformed
normalizers,
\[
  \widehat\ell_T^{\rm sv}(\theta;B)
  =
  \sum_{t=1}^T \log \widehat Z_t^{\rm sv}(\theta;B),
  \qquad
  \widehat Z_t^{\rm sv}(\theta;B)\approx Z_t(\theta).
\]
Only after that declaration does the same-scalar derivative contract become
meaningful.  Direct-likelihood quadrature, fixed-cloud reweighting, and
Zhao--Cui comparisons are legitimate only when they approximate this same
normalizer sequence.  The previous actual-SV ``Lane~B'' framing violated that
ordering by differentiating a Gaussian-closure replacement scalar built from
predictive observation moments of an augmented variable.  That scalar was not a
second admissible target for the intended inference problem; it was the wrong
scalar.  A finite-difference or autodiff check on that wrong scalar can show only
internal consistency of the replacement construction, not correctness for the
actual-SV likelihood.

\section{Finite-Difference Contract And Branch-Validity Logic}
\label{sec:bf-hd-fixed-branch-fd-contract}

The finite-difference contract is
\begin{equation}
  D_i(h)=
  \frac{\widehat\ell_T(\theta_0+h e_i;B)-\widehat\ell_T(\theta_0-h e_i;B)}{2h},
  \label{eq:bf-hd-fixed-branch-fd}
\end{equation}
with both perturbed evaluations reusing the same fixed ledger while recomputing
all parameter-dependent numerical values through the same declared equations.

For the retained-object lane, this means the domains, points, ranks, shifts, and
sweep order are unchanged, but the fitted core values
\begin{equation}
  C_{t,k}(\theta_0+h e_i;B)
  \quad\text{and}\quad
  C_{t,k}(\theta_0-h e_i;B)
  \label{eq:bf-hd-ttkr-perturbed-cores}
\end{equation}
are recomputed from the same saved fitting rule.  They are not copied from the
base parameter.  A row is branch-valid only if the two perturbations share:
\begin{enumerate}[label=(\roman*), leftmargin=0.28in]
  \item the same saved fixed ledger, and
  \item the same realized acceptance/failure path through the declared branch
  logic.
\end{enumerate}
If either side changes the accepted/failure pattern, the row is branch-invalid
and is not evidence about the derivative of the same scalar.

\paragraph{Copied-core and branch-identity audit protocol.}
The operational diagnostic should promote these prose rules into explicit
indicators.  For a step ladder \(h_i\), define
\begin{equation}
  I_\pm(h_i)=\mathbf{1}\{\mathcal H_{\rm fixed}(\theta_0\pm h_i e_j)=\mathcal H_{\rm fixed}(\theta_0)\},
  \label{eq:bf-hd-ttkr-branch-indicator}
\end{equation}
which checks branch-identity equality of the frozen fields, and
\begin{equation}
  K_\pm(h_i)=\mathbf{1}\{C_{t,k}(\theta_0\pm h_i e_j;B)\ \text{was recomputed by the fixed fitting rule}\},
  \label{eq:bf-hd-ttkr-copied-core-indicator}
\end{equation}
which checks that perturbed cores were not copied from the base parameter.
Record a decreasing-window indicator
\begin{equation}
  W_i=\mathbf{1}\{e(h_i)>e(h_{i+1})>e(h_{i+2})\},
  \label{eq:bf-hd-ttkr-decreasing-window}
\end{equation}
and classify the diagnostic status as one of:
\begin{equation}
\text{status}=
\begin{cases}
\text{branch identity failure},&\min_{i,\pm} I_\pm(h_i)=0,\\
\text{copied-core failure},&\min_{i,\pm} K_\pm(h_i)=0,\\
\text{decreasing-window agreement},&\max_i W_i=1,\\
\text{inconclusive: no decreasing window},&\text{otherwise}.
\end{cases}
  \label{eq:bf-hd-ttkr-fd-status}
\end{equation}
This makes the same-branch contract auditable.  If a frozen field differs, the
centered difference is comparing different functions.  If a differentiable field
is copied rather than recomputed, the centered difference is not evaluating the
fixed fitting rule.

\section{Value Path, Derivative Path, And What Is Actually Differentiated}
\label{sec:bf-hd-derivative-ledger}

The implementation burden of a same-scalar chapter is not just to state that a
gradient exists.  It must say which objects belong to the value path and which
objects belong to the derivative path.  A generic fixed-branch workflow has the
shape
\begin{equation}
\begin{aligned}
  &\text{value path:}
    &&\text{frozen structure} \to \widehat Z_t \to \widehat\ell_T,\\
  &\text{derivative path:}
    &&\text{same frozen structure} \to \dot{\widehat Z}_t \to
      \partial_i\widehat\ell_T.
\end{aligned}
  \label{eq:bf-hd-value-derivative-ledger}
\end{equation}
For the sparse-grid lane this means differentiating the same fixed cloud and the
same branch-conditioned Gaussian value path.  For the retained-object lane it
means differentiating the same retained-object fitting equations, mass
contractions, quotient updates, and carried-filter update equations without
changing domains, points, ranks, shifts, or sweep order.

The most important implementation distinction is between:
\begin{itemize}[leftmargin=0.28in]
  \item objects needed to close the current scalar derivative, and
  \item objects needed only to propagate the branch-consistent state to the next
  step.
\end{itemize}
This distinction is what keeps a same-scalar derivative from being confused with
a numerically similar but structurally different adaptive algorithm.

\paragraph{Retained-object derivative mechanics on the frozen branch.}
For the TT/KR lane, the forward and derivative paths should be read as the same
finite sequence of fixed equations with different outputs.  A useful teaching
layer is:
\begin{equation}
\begin{array}{p{0.41\linewidth}p{0.45\linewidth}}
\toprule
Forward scalar path & Derivative path\\
\midrule
\(\widetilde q_{t,j}=\widehat p_{t-1,j} f_j g_j J_j\) &
\(\dot{\widetilde q}_{t,j}=J_j(\dot{\widehat p}_{t-1,j}f_jg_j+\widehat p_{t-1,j}\dot f_jg_j+\widehat p_{t-1,j}f_j\dot g_j)\)\\
\(y_j=e^{c_t/2}\sqrt{\widetilde q_{t,j}}\) &
\(\dot y_j=e^{c_t/2}\dot{\widetilde q}_{t,j}/(2\sqrt{\widetilde q}_{t,j})\)\\
\(A_k\to N_k=A_k^\top W A_k+\rho I,\ d_k=A_k^\top Wy\) &
\(\dot A_k\to \dot N_k=\dot A_k^\top WA_k+A_k^\top W\dot A_k,\ \dot d_k=\dot A_k^\top Wy+A_k^\top W\dot y\)\\
\(N_k g_k=d_k\) &
\(N_k\dot g_k=\dot d_k-\dot N_k g_k\)\\
\(g_k\to C_k\to M_{>k}\to \widehat Z_t\) &
\(\dot g_k\to \dot C_k\to \dot M_{>k}\to \dot{\widehat Z}_t\)\\
\(\widehat p_t=a_t/\widehat Z_t\) &
\(\dot{\widehat p}_t=\dot a_t/\widehat Z_t-a_t\dot{\widehat Z}_t/\widehat Z_t^2\)\\
\bottomrule
\end{array}
  \label{eq:bf-hd-ttkr-derivative-teaching-table}
\end{equation}
The fixed linear solve is the key finite-dimensional derivative step:
\begin{equation}
  N_k g_k=d_k
  \quad\Longrightarrow\quad
  N_k\dot g_k=\dot d_k-\dot N_k g_k.
  \label{eq:bf-hd-ttkr-fixed-solve-derivative}
\end{equation}
Computation should solve \eqref{eq:bf-hd-ttkr-fixed-solve-derivative} with the
same numerical linear algebra used in the forward pass; it should not form an
explicit inverse.  After the core derivative is obtained, the remaining objects
follow by repeated product-rule and quotient-rule propagation through the same
frozen branch.

\paragraph{Mass-contraction and normalizer block.}
The squared-TT normalizer is not a separate conceptual object from the fitted
cores; it is a contraction of those same cores against the frozen basis mass
matrices.  Let the right mass contraction satisfy a recursion of the form
\begin{equation}
  M_{>j-1}[a,a']=
  \sum_{b,b',\ell,\ell'}
  C_j[a,\ell,b]
  M_{>j}[b,b']
  C_j[a',\ell',b']
  B_j[\ell,\ell'],
  \label{eq:bf-hd-ttkr-mass-recursion}
\end{equation}
with basis mass matrix \(B_j\) frozen by the branch.  Differentiate it by the
product rule:
\begin{equation}
\begin{aligned}
  \dot M_{>j-1}[a,a']
  &=
  \sum_{b,b',\ell,\ell'}
  \dot C_j[a,\ell,b]
  M_{>j}[b,b']
  C_j[a',\ell',b']
  B_j[\ell,\ell']\\
  &\quad+
  \sum_{b,b',\ell,\ell'}
  C_j[a,\ell,b]
  \dot M_{>j}[b,b']
  C_j[a',\ell',b']
  B_j[\ell,\ell']\\
  &\quad+
  \sum_{b,b',\ell,\ell'}
  C_j[a,\ell,b]
  M_{>j}[b,b']
  \dot C_j[a',\ell',b']
  B_j[\ell,\ell'].
\end{aligned}
  \label{eq:bf-hd-ttkr-dot-mass-recursion}
\end{equation}
Starting from the terminal contraction, this recursion produces the square mass
\(R_t\) and its derivative \(\dot R_t\).  If
\begin{equation}
  \zeta_t=R_t+\tau_t,
  \qquad
  \widehat Z_t=e^{-c_t}\zeta_t,
  \label{eq:bf-hd-ttkr-normalizer-block}
\end{equation}
then, with \(c_t\) and \(\tau_t\) frozen on the branch,
\begin{equation}
  \dot{\widehat Z}_t=e^{-c_t}\dot R_t.
  \label{eq:bf-hd-ttkr-dot-normalizer}
\end{equation}
This is the mass-contraction block the derivative must preserve.  Omitting one of
the product-rule terms in \eqref{eq:bf-hd-ttkr-dot-mass-recursion} means the
reported gradient no longer differentiates the same scalar normalizer.

\paragraph{Carried-filter and next-step closure block.}
The retained numerator is another contraction of the same represented
adjacent-state object.  Write
\begin{equation}
  a_t(z_t;\beta)=
  \int \widehat q_t(z_t,z_{t-1};\beta)\dd z_{t-1},
  \qquad
  \dot a_t(z_t;\beta)=
  \int \dot{\widehat q}_t(z_t,z_{t-1};\beta)\dd z_{t-1}.
  \label{eq:bf-hd-ttkr-carried-numerator}
\end{equation}
The normalized carried filter then satisfies the quotient rule
\begin{equation}
  \dot{\widehat p}_t(z_t;\beta)=
  \frac{\dot a_t(z_t;\beta)}{\widehat Z_t(\beta)}
  -
  \frac{a_t(z_t;\beta)\dot{\widehat Z}_t(\beta)}{\widehat Z_t(\beta)^2}.
  \label{eq:bf-hd-ttkr-dot-carried-filter}
\end{equation}
This is the crucial bridge between time steps.  At time \(t+1\), the fixed-branch
target contains \(\widehat p_t\), so its derivative contains
\(\dot{\widehat p}_t\).  Without saving the carried-filter derivative, the
next-step gradient is incomplete even if the current-step score is numerically
correct.

\paragraph{Saved retained-filter object and query rule.}
The phrase “save the carried object” must denote a finite-dimensional evaluator,
not an informal function name.  For retained state dimension $m>1$, let the
retained block use the product basis
\begin{equation}
  b_t(z_t)=b_{t,1}(z_{t,1})\otimes\cdots\otimes b_{t,m}(z_{t,m})
  \in\R^{p_{\rm ret}},
  \qquad
  p_{\rm ret}=\prod_{j=1}^m p_j.
  \label{eq:bf-hd-ttkr-retained-product-basis}
\end{equation}
The retained numerator is represented by a coefficient matrix
\begin{equation}
  a_t(z_t;\beta)=b_t(z_t)^\top Q_t(\beta)b_t(z_t),
  \qquad
  Q_t\in\R^{p_{\rm ret}\times p_{\rm ret}},
  \label{eq:bf-hd-ttkr-retained-Q}
\end{equation}
and its derivative uses the matching dotted matrix
\begin{equation}
  \dot a_t(z_t;\beta)=b_t(z_t)^\top \dot Q_t(\beta)b_t(z_t),
  \qquad
  \dot Q_t\in\R^{p_{\rm ret}\times p_{\rm ret}}.
  \label{eq:bf-hd-ttkr-retained-dot-Q}
\end{equation}
Normalize the stored evaluator by
\begin{equation}
  P_t=\frac{Q_t}{\widehat Z_t},
  \qquad
  \dot P_t=
  \frac{\dot Q_t}{\widehat Z_t}
  -
  \frac{Q_t\dot{\widehat Z}_t}{\widehat Z_t^2}.
  \label{eq:bf-hd-ttkr-retained-P}
\end{equation}
For next-step query points $z_t^{\rm query}(j)$, build the retained-basis matrix
\begin{equation}
  B^{\rm query}[j,:]=b_t(z_t^{\rm query}(j))^\top,
  \qquad
  B^{\rm query}\in\R^{M\times p_{\rm ret}},
  \label{eq:bf-hd-ttkr-retained-query-basis}
\end{equation}
and evaluate the saved carried object by
\begin{equation}
  \widehat p_t^{\rm ref}[j]
  =
  B^{\rm query}[j,:]P_tB^{\rm query}[j,:]^\top,
  \qquad
  \dot{\widehat p}_t^{\rm ref}[j]
  =
  B^{\rm query}[j,:]\dot P_tB^{\rm query}[j,:]^\top.
  \label{eq:bf-hd-ttkr-retained-query-rule}
\end{equation}
The next-step query rule is that $z_t^{\rm query}(j)$ is the previous-state block
of the fixed branch's fitting point at time $t+1$.  A later implementation may
store $Q_t$ in TT, low-rank, or sparse form, but only if it preserves the same
evaluator and derivative evaluator in
\eqref{eq:bf-hd-ttkr-retained-query-rule}.  This is the stored-field convention
that closes the derivative recursion on the same branch.

\paragraph{Branch-manifest reporting hook.}
Derivative and finite-difference reports for this lane should therefore carry a
TT/KR branch-manifest identifier together with the saved-field convention, so
that branch-valid rows certify equality of the actual retained evaluator rather
than only superficial notation.

\paragraph{Compact derivative-reading checklists.}
A later implementation review should be able to verify the derivative lane block
by block:
\begin{equation}
\begin{array}{ll}
\text{squared-density block:} & \overline q_t=e^{-c_t}(\phi_t^2+\tau_t\lambda_t),\ \partial(\phi_t^2)=2\phi_t\dot\phi_t;\\[1mm]
\text{mass-contraction block:} & R_t=\int \phi_t^2,\ \dot R_t\ \text{from product-rule contractions};\\[1mm]
\text{fixed-solve block:} & N_kg_k=d_k\ \mapsto\ N_k\dot g_k=\dot d_k-\dot N_kg_k;\\[1mm]
\text{carried-filter block:} & \widehat p_t=a_t/\widehat Z_t,\ \dot{\widehat p}_t=\dot a_t/\widehat Z_t-a_t\dot{\widehat Z}_t/\widehat Z_t^2.
\end{array}
  \label{eq:bf-hd-ttkr-derivative-checklists}
\end{equation}
The corresponding failure-mode summaries are:
\begin{itemize}[leftmargin=0.28in]
  \item \textbf{squared-density failure:} differentiating a different fitted
  square root than the one actually squared,
  \item \textbf{mass-contraction failure:} omitting a product-rule term in the
  contraction recursion,
  \item \textbf{fixed-solve failure:} changing the normal equations or copying a
  stale core at perturbed parameters,
  \item \textbf{carried-filter failure:} feeding time \(t+1\) without the saved
  dotted retained object or with an undeclared query rule.
\end{itemize}
These lists are not replacement derivations; they are compact audit objects that
make it harder for a future implementation to silently drift away from the
mathematical branch defined here.

\section{Approximate Targets, Exact-Target Corrections, And HMC Admissibility}
\label{sec:bf-hd-fixed-branch-hmc}

For an unconstrained coordinate \(q\in\R^d\) and transform
\(\theta=\tau(q)\), the exact-target HMC potential is
\begin{equation}
  U(q)
  =
  -\ell_T(\tau(q))-
  \log \pi(\tau(q))-
  \log\left|\det D\tau(q)\right|.
  \label{eq:bf-hd-hmc-potential-live}
\end{equation}
If the filter supplies only \(\widehat\ell_T\), then ordinary HMC samples the
declared approximate posterior defined by \(\widehat\ell_T\).  Recovering the
exact posterior from an approximate or randomized likelihood would require a
separately defined exact-target extended-state or Metropolis correction scheme;
such a scheme is outside the scope of this chapter.

The HMC admissibility burden is therefore downstream of same-scalar discipline:
\begin{enumerate}[label=(\arabic*), leftmargin=0.28in]
  \item the filter must define the scalar likelihood or approximate likelihood
  before any sampler can be interpreted,
  \item the gradient used by HMC must be the gradient of the same scalar used in
  the Metropolis correction and posterior interpretation,
  \item transformed targets, including TT/KR-style maps, require the
  Jacobian-corrected target and valid support.
\end{enumerate}

\section{Exports To Validation And Promotion Chapters}
\label{sec:bf-hd-fixed-branch-exports}

The later validation chapter may import only the following consequences from the
present chapter:
\begin{itemize}[leftmargin=0.28in]
  \item the branch-valid finite-difference contract,
  \item the distinction between fixed structural ledgers and recomputed
  differentiable ledgers,
  \item the exact-target versus approximate-target HMC admissibility boundary,
  \item the rule that divergences, nonfinite gradients, and same-scalar mismatch
  are veto diagnostics rather than efficiency datapoints.
\end{itemize}

\section{Implementation Boundary}
\label{sec:bf-hd-hmc-boundary}

BayesFilter has no nonlinear HMC convergence result in this lane, no NeuTra
backend, no RMHMC backend, no tensor/KR transport-preconditioned HMC backend,
and no NAWM posterior recovery benchmark.  This chapter defines a shared branch,
scalar, and derivative contract.  It does not promote any sampler or backend as
production-ready for nonlinear state-space models.
